\documentclass[11pt,a4paper]{article}
\usepackage[utf8]{inputenc}
\usepackage[T1]{fontenc}
\usepackage[spanish,es-noquoting]{babel}
\usepackage{amsmath,amssymb}
\usepackage{graphicx}
\usepackage{booktabs}
\usepackage{siunitx}
\usepackage[margin=2.5cm]{geometry}
\usepackage{hyperref}
\usepackage{xcolor}
\hypersetup{colorlinks=true,linkcolor=blue!50!black,urlcolor=blue!50!black}
\sisetup{group-separator={,},group-minimum-digits=4}

\title{Aceleraci\'on y paralelizaci\'on de simulaciones Monte Carlo\\
       de nanodiscos magn\'eticos}
\author{Estudio de simulaci\'on --- Parte 1 del plan de pruebas}
\date{\today}

\begin{document}
\maketitle

\begin{abstract}
Se mide qu\'e aporta cada mecanismo de aceleraci\'on en una simulaci\'on Monte
Carlo de un nanodisco magn\'etico, y d\'onde conviene poner el paralelismo una
vez que el algoritmo es correcto. Las dos aceleraciones del m\'etodo ---
descomposici\'on matricial (MD) y actualizaci\'on bipartita (BP) --- se ablacionan
por separado sobre CPU y sobre 1, 2 y 4 GPU H100. Se encuentra que BP no es un
acelerador sino un requisito de correctitud que \emph{cuesta} un factor 2 en la
implementaci\'on vectorizada, y que MD solo acelera \emph{porque} BP vuelve legal
la actualizaci\'on simult\'anea: aplicada a una cadena secuencial la empeora 348
veces. Una GPU sin MD resulta 438 veces m\'as lenta que una CPU. Sobre el eje de
paralelizaci\'on se comparan tres agrupaciones del mismo trabajo ---sobre
r\'eplicas, sobre estados, o ambas--- y se establece que el n\'umero de
dispositivos es gratuito a lote fijo (0.9\,\% entre 1 y 4 GPU) mientras que el
lote por dispositivo tiene un \'optimo cercano a 50 cadenas con 27\,\% de castigo
en los extremos. Finalmente se verifica que ninguna de estas decisiones altera la
f\'isica: seis paralelizaciones distintas concuerdan en la energ\'ia al
0.07\,\%, aunque produzcan realizaciones espacialmente no correlacionadas.
\end{abstract}

\tableofcontents

\section{El sistema y el m\'etodo}

\subsection{Hamiltoniano}

Se simula un nanodisco magn\'etico sobre una red c\'ubica simple, con
$R_d = 18.30$ en unidades de red, $L = 5$ capas y $N = 39$ sitios por lado, lo
que deja \num{5245} sitios activos dentro del disco. El Hamiltoniano extendido es
\begin{equation}
\mathcal{H} = -J_{\mathrm{ex}}\!\!\sum_{\langle ij\rangle}\! \mathbf{S}_i\cdot\mathbf{S}_j
 \; - \; \mathcal{H}_{\mathrm{DM}} \; + \; \mathcal{H}_{\mathrm{an}}
 \; - \; H_{\mathrm{ex}}\!\sum_i \left(S_i^x\cos\gamma + S_i^z\sin\gamma\right),
\end{equation}
con intercambio a primeros vecinos, interacci\'on Dzyaloshinskii--Moriya
interfacial, anisotrop\'ia c\'ubica y t\'ermino Zeeman.

\paragraph{DMI interfacial.} El t\'ermino DM act\'ua \textbf{solo sobre los
enlaces en el plano} ($x$ e $y$) y con signo negativo:
\begin{equation}
\mathcal{H}_{\mathrm{DM}} = K_{\mathrm{DM}}\sum_i
 \Big[ S_i^y S_{i+\hat x}^z - S_i^z S_{i+\hat x}^y
     + S_i^z S_{i+\hat y}^x - S_i^x S_{i+\hat y}^z \Big].
\end{equation}
Esto no es una omisi\'on: en una pel\'icula delgada el vector DM proviene de la
interfaz y act\'ua sobre enlaces en el plano. Incluir el enlace en $z$
modelar\'ia DMI de volumen, que es otra clase de material. La diferencia es
medible y grande: sobre la misma configuraci\'on de espines, usar los tres ejes
con signo $+$ cambia la energ\'ia total en un \textbf{50\,\%}.

\paragraph{Anisotrop\'ia c\'ubica.} $K_1$ es el coeficiente de primer orden y
$K_2$ el de segundo, multiplicando combinaciones de los cosenos directores
$\alpha_i = S_i/S_o$ que definen la direcci\'on del momento:
\begin{equation}
E_{\mathrm{an}} = K_1\left(\alpha_1^2\alpha_2^2 + \alpha_1^2\alpha_3^2
 + \alpha_2^2\alpha_3^2\right) + K_2\,\alpha_1^2\alpha_2^2\alpha_3^2 .
\end{equation}
La implementaci\'on usa $K_1/S_o^4$ y $K_2/S_o^6$, que es exactamente lo que
convierte componentes de esp\'in en cosenos directores; la equivalencia se
verific\'o num\'ericamente a precisi\'on de m\'aquina. $K_1$ toma el valor de
superficie $K_{\mathrm{anS}}$ en los sitios con menos de seis vecinos dentro del
disco.

\subsection{Protocolo}

Recocido simulado desde \SI{20.0}{\kelvin} hasta \SI{0.1}{\kelvin} en
\textbf{200 puntos de temperatura}, con \num{10000} barridos de termalizaci\'on y
\num{5000} de medici\'on en cada uno: \num{3000000} barridos por escalera
completa. Un \textbf{barrido} (equivalentemente, un paso Monte Carlo, MCS) es un
intento de actualizaci\'on por cada uno de los \num{5245} espines.

\subsection{Una trampa num\'erica: la varianza en \texttt{float32}}

$C_v$ y $\chi$ son varianzas, y acumular $E$ y $E^2$ directamente las destruye en
precisi\'on simple. A \SI{0.2}{\kelvin} una corrida tiene
$\langle E\rangle \approx -1.7\times10^4$, de modo que
$\langle E^2\rangle \approx 2.9\times10^8$, donde un ulp de \texttt{float32} vale
unos 34 --- mayor que la varianza misma. La diferencia
$\langle E^2\rangle - \langle E\rangle^2$ sale \textbf{negativa} y $C_v$ marca
$-437$. La soluci\'on es acumular la fluctuaci\'on respecto de una referencia
tomada una vez por temperatura: un corrimiento constante deja la varianza
intacta, as\'i que es exacto.

\section{Las dos aceleraciones}

\subsection{Por qu\'e es necesaria la bipartidad}

Cada esp\'in decide si cambia mirando a sus vecinos. Si dos vecinos deciden
simult\'aneamente, cada uno observa el valor \emph{viejo} del otro: el cambio de
energ\'ia conjunto no es la suma de los individuales y la cadena deja de muestrear
la distribuci\'on de Boltzmann.

La soluci\'on es colorear la red con la paridad $(x+y+z)\bmod 2$. Ning\'un sitio
par toca a otro par, de modo que \textbf{todos los pares pueden decidir a la vez}
sin inconsistencia. Un barrido son entonces \textbf{dos pasadas}, una por color,
y cada esp\'in recibe exactamente una oportunidad: los mismos \num{5245} intentos.

Que esto sea un reordenamiento \emph{exacto} del Metropolis secuencial no se
asume: se demuestra. Alimentando ambos lazos con propuestas id\'enticas y
uniformes id\'enticas, el resultado coincide bit a bit; y a lo largo de una
escalera completa de 200 puntos la energ\'ia por esp\'in del brazo secuencial por
color coincide con la del tablero vectorizado con desviaci\'on
\textbf{m\'axima de 0.000\,\%}.

\paragraph{Lo que cuesta.} La bipartidad \textbf{no es gratis}, y la medici\'on
depende de la implementaci\'on:

\begin{table}[h]\centering
\begin{tabular}{lrrr}
\toprule
implementaci\'on & con BP & sin BP & costo de BP \\
\midrule
escalar, sitio por sitio & 12.77 & 13.03 & \textbf{0.98$\times$} \\
numpy vectorizado        & 156.69 & 66.68 & \textbf{2.35$\times$} \\
JAX (el n\'ucleo que se usa) & 65.16 & 32.23 & \textbf{2.02$\times$} \\
\bottomrule
\end{tabular}
\caption{Costo de la actualizaci\'on bipartita, en ms por barrido a 100 cadenas.
La versi\'on ``sin BP'' actualiza todo el disco en una sola pasada: es el doble de
r\'apida y muestrea la distribuci\'on equivocada.}
\end{table}

En una implementaci\'on escalar el coloreado es gratis porque solo cambia el
orden de visita. En una vectorizada cuesta un factor 2, y la raz\'on est\'a en el
lazo: cada pasada calcula $\Delta E$ sobre \emph{toda} la red y descarta la mitad
que no le corresponde. Mismos \num{5245} intentos, el doble de aritm\'etica.

\paragraph{Qu\'e se rompe sin ella.} La versi\'on de una pasada no es una cadena
m\'as r\'apida, es \emph{otra} cadena. La energ\'ia se desplaza menos de 1\,\% y
sola habr\'ia ocultado el problema; lo que se mueve es $|M|$ entre 2 y 5\,\% y la
carga topol\'ogica de Berg--L\"uscher de la capa central \textbf{hasta una unidad
completa}. Esa es precisamente la cantidad con la que est\'an etiquetados los
estados del conjunto de datos liberado.

\subsection{Por qu\'e es necesaria la descomposici\'on matricial}

La descomposici\'on matricial expresa la energ\'ia de prueba de muchos sitios
como operaciones sobre arreglos en lugar de aritm\'etica sitio por sitio. Es lo
que \textbf{crea el ancho} de trabajo paralelo.

\begin{table}[h]\centering
\begin{tabular}{lrrr}
\toprule
m\'etodo & pasos en serie & cuentas & ms/barrido \\
\midrule
sitio por sitio        & \num{5245} & $1\times$ & 352.3 \\
con BP (dos pasadas)   & 2          & $2\times$ & \textbf{65.0} \\
sin BP (una pasada)    & 1          & $1\times$ & 31.9 \\
\bottomrule
\end{tabular}
\caption{Misma m\'aquina, 100 cadenas. La fila de una pasada es incorrecta y se
incluye solo como cota.}
\end{table}

El resultado parece contradictorio ---el m\'etodo con \emph{m\'as} cuentas es
cinco veces m\'as r\'apido--- y no lo es: lo que cuesta tiempo no es la cantidad
de operaciones sino cu\'antas veces hay que detenerse a esperar un resultado.
Hacer el doble de aritm\'etica en 2 esperas le gana por lejos a hacer la mitad en
\num{5245}.

\paragraph{MD y BP no son separables.} Medido a 100 cadenas sobre CPU:
\begin{itemize}
  \item BP sola: $0.98\times$ --- no acelera en implementaci\'on escalar.
  \item MD con BP presente: $\mathbf{5.87\times}$.
  \item MD \emph{sin} BP: $0.0029\times$, es decir \textbf{348 veces m\'as lenta}.
\end{itemize}
Aplicar la descomposici\'on a una cadena secuencial es una \emph{pesimizaci\'on}:
construye el campo de vecinos de toda la red y descarta todo menos un elemento.
MD acelera \emph{porque} BP vuelve legal la actualizaci\'on simult\'anea.

\paragraph{Una GPU sin MD es peor que no tener GPU.} Sin descomposici\'on la
cadena es secuencial y cada sitio se convierte en un despacho independiente al
dispositivo: \num{5245} despachos por barrido, cada uno con su latencia de
lanzamiento y sincronizaci\'on, para calcular unas pocas docenas de flotantes.
Medido, una H100 corre a \SI{146315}{\milli\second} por barrido contra
\SI{334}{\milli\second} de la CPU: \textbf{438 veces m\'as lenta}. Y empeora al
agregar dispositivos, porque cada despacho debe alcanzar y sincronizar $N$ de
ellos.

\begin{figure}[h]\centering
\includegraphics[width=\textwidth]{figures/ablation_grid.pdf}
\caption{Ablaci\'on MD $\times$ BP sobre cuatro configuraciones de hardware, a
100 cadenas y geometr\'ia de producci\'on. Escala logar\'itmica. Solo la celda
MD\,+\,BP se beneficia de una GPU; las otras tres corren m\'as lento en GPU que en
CPU, y peor cuantas m\'as GPU.}
\end{figure}

\section{D\'onde poner el paralelismo}

Fijado el algoritmo, queda decidir c\'omo repartir el trabajo. Se comparan tres
agrupaciones de un mismo total de \textbf{4 estados $\times$ 100 r\'eplicas = 400
cadenas independientes}, cada una recorriendo la escalera completa de 200 puntos.

\begin{table}[h]\centering
\begin{tabular}{llp{6.0cm}}
\toprule
modo & agrupaci\'on & eje que usa \\
\midrule
BP solo                & 400 lanzamientos de 1 cadena   & solo la actualizaci\'on bipartita \\
BP + r\'eplicas        & 4 lanzamientos de 100 cadenas  & lote sobre semillas, un estado por vez \\
BP + estados           & 100 lanzamientos de 4 cadenas  & lote sobre par\'ametros, una semilla por vez \\
BP + r\'epl. + estados & 1 lanzamiento de 400 cadenas & ambos ejes \\
\bottomrule
\end{tabular}
\end{table}

Los cuatro estados son las medianas de cuatro c\'umulos de fase del conjunto de
datos liberado ---\texttt{uniform}, \texttt{helical}, \texttt{points}
(skyrmiones) y \texttt{lines} (laberinto)--- de modo que son configuraciones que
existen y no esquinas inventadas. Los cuatro tienen $J_2=J_3=J_4=0$, lo que hace
que el coloreado bipartito sea exacto para todos; el c\'odigo lo verifica con una
aserci\'on en lugar de confiar.

\subsection{Resultados}

\begin{table}[h]\centering
\begin{tabular}{lrrrr}
\toprule
modo & CPU & 1$\times$H100 & 2$\times$H100 & 4$\times$H100 \\
\midrule
BP solo                   & 79.4 d$^*$ & 5.8 h$^*$ & 5.8 h$^*$ & 5.8 h$^*$ \\
\textbf{BP + r\'eplicas}  & 10.3 d$^*$ & \textbf{51.7 min} & \textbf{27.6 min} & \textbf{16.9 min} \\
BP + estados              & 45.4 d$^*$ & 1.8 h$^*$ & 1.8 h$^*$ & 1.7 h$^*$ \\
\textbf{BP + r\'epl. + est.} & 7.6 d$^*$ & \textbf{1.0 h} & \textbf{30.1 min} & \textbf{16.4 min} \\
\midrule
\multicolumn{5}{l}{\small escalado $1\to4$ dispositivos:}\\
BP solo & \multicolumn{4}{l}{$0.99\times$} \\
BP + r\'eplicas & \multicolumn{4}{l}{$3.06\times$} \\
BP + estados & \multicolumn{4}{l}{$1.06\times$} \\
BP + r\'epl. + est. & \multicolumn{4}{l}{$3.80\times$ \quad (ideal $4.00\times$)} \\
\bottomrule
\end{tabular}
\caption{Tiempo total de ejecuci\'on. Las celdas sin asterisco son escaleras
completas medidas de punta a punta; las marcadas con $^*$ se extrapolan de un
ritmo por barrido medido.}
\end{table}

\begin{figure}[h]\centering
\includegraphics[width=\textwidth]{figures/total_time_all.pdf}\\[1ex]
\includegraphics[width=\textwidth]{figures/total_time_gpu.pdf}
\caption{Tiempo total de ejecuci\'on por dispositivo. Arriba, escala
logar\'itmica incluyendo CPU; abajo, solo GPU en escala lineal. Las barras
rayadas son extrapolaciones.}
\end{figure}

\subsection{Por qu\'e paralelizar por experimento casi no ayuda}

\textbf{$1.06\times$ con cuatro GPU.} La agrupaci\'on por estados lanza un lote de
\textbf{4 cadenas}, cien veces. Cuatro cadenas no se acercan a llenar una H100
---la saturaci\'on llega alrededor de 25--- y agregar dispositivos lo
\emph{empeora}, porque esas cuatro cadenas se reparten en una por dispositivo: se
pasa de una GPU casi ociosa a cuatro completamente ociosas.

La penalizaci\'on contra el mejor modo \emph{crece} con el hardware: $2.07\times$
con una GPU, $3.83\times$ con dos, $6.19\times$ con cuatro. Es lo contrario de
para lo que se compra hardware.

Esto importa de forma directa porque \textbf{es la agrupaci\'on que el instinto
elige para generar un conjunto de datos}: un conjunto de datos var\'ia el estado,
no la semilla. La recomendaci\'on es juntar \emph{muchos} estados en un lote
ancho; al dispositivo le da casi igual de d\'onde viene el ancho, con tal de que
lo haya.

\subsection{Por qu\'e el paralelismo por r\'eplicas rinde bien}

Las r\'eplicas son cadenas de Markov \emph{independientes}: no hay ninguna
operaci\'on colectiva entre ellas, nada que sincronizar hasta el final. El eje de
r\'eplicas produce lotes de 100 cadenas, muy por encima de la saturaci\'on, y
escala $3.06\times$ sobre cuatro dispositivos. Agrupar ambos ejes a la vez da
$3.80\times$, cerca del ideal.

Entre esas dos la elecci\'on apenas importa: 16.9 contra 16.4 minutos sobre cuatro
GPU. \textbf{Lo que importa es tener un lote ancho}, no de qu\'e eje provenga.

\subsection{Lote por dispositivo contra n\'umero de dispositivos}

Las seis escaleras suger\'ian un costo por multi-dispositivo. Dos barridos, cada
uno dentro de un mismo contenedor y con tres repeticiones por celda, lo separan.
Se descarta siempre la primera repetici\'on: carga la compilaci\'on JIT.

\begin{table}[h]\centering
\begin{tabular}{rr@{\qquad}rrr}
\toprule
\multicolumn{2}{c}{lote fijo $B=25$} & \multicolumn{3}{c}{1 GPU fija} \\
GPU & $\mu$s/cadena-barrido & $B$ & $\mu$s/cadena-barrido & sd \\
\midrule
1 & 2.826 & 25  & 2.826 & 0.0010 \\
2 & 2.800 & \textbf{50}  & \textbf{2.323} & 0.0020 \\
4 & 2.807 & 100 & 2.960 & 0.0020 \\
\midrule
\multicolumn{2}{l}{dispersi\'on \textbf{0.9\,\%}} & \multicolumn{3}{l}{castigo \textbf{27\,\%} en los extremos} \\
\bottomrule
\end{tabular}
\caption{Tiempo-dispositivo por cadena-barrido, escaleras completas.}
\end{table}

\begin{figure}[h]\centering
\includegraphics[width=\textwidth]{figures/batch_device_scan.pdf}
\caption{Izquierda: efecto puro del lote por dispositivo, con el n\'umero de GPU
fijo en uno. Derecha: efecto puro del n\'umero de dispositivos, con el lote fijo
en 25 cadenas por GPU.}
\end{figure}

\paragraph{Coordinar dispositivos es gratuito.} A lote fijo, pasar de una a
cuatro GPU cuesta 0.9\,\%. Las r\'eplicas no tienen colectivo y \texttt{pmap} no
agrega sobrecosto medible.

\paragraph{El lote por dispositivo s\'i importa.} Hay un m\'inimo claro cerca de
$B=50$ y un castigo del 27\,\% en los extremos, con desviaciones de 0.002: es
se\~nal. La recomendaci\'on pr\'actica es \textbf{apuntar a unas 50 cadenas por
GPU}, y explica por qu\'e la agrupaci\'on por r\'eplicas sobre cuatro
dispositivos ($B=25$) rinde por debajo de su propio escalado.

\paragraph{Lo que no queda aislado.} El barrido no reproduce los valores
absolutos de las escaleras: a $B=100$ la escalera mide 2.586 y el barrido 2.960,
un 14\,\% de diferencia. Queda un tercer factor sin aislar; dos candidatos no
probados son que las escaleras pasan par\'ametros como arreglos por cadena
mientras el barrido usa escalares, y que las escaleras cronometran la primera
llamada con la compilaci\'on incluida. La \emph{forma} de la curva y el resultado
sobre dispositivos quedan establecidos; los valores absolutos del barrido no
deben citarse como costo de escalera.

\section{En qu\'e afecta la paralelizaci\'on a la formaci\'on de estructuras}

Las seis corridas difieren \emph{\'unicamente} en c\'omo se entregaron las mismas
400 cadenas al hardware. El Hamiltoniano, la actualizaci\'on bipartita, la
escalera de temperaturas y el n\'umero de r\'eplicas son id\'enticos.

\subsection{Los observables coinciden}

\begin{table}[h]\centering
\begin{tabular}{lrrrrr}
\toprule
estado & $E$ & $M$ & $C_v$ & $\chi$ & $Q$ \\
\midrule
uniform & 0.02\,\% & 0.02\,\% & 6.1\,\% & 3.3\,\% & 5.1\,\% \\
helical & 0.07\,\% & 1.8\,\%  & 8.2\,\% & 13.2\,\% & 15.8\,\% \\
points  & 0.04\,\% & 0.29\,\% & 7.2\,\% & 18.7\,\% & 2.5\,\% \\
lines   & 0.05\,\% & 9.5\,\%  & 8.0\,\% & 25.9\,\% & 34.7\,\% \\
\bottomrule
\end{tabular}
\caption{Dispersi\'on m\'axima entre las seis paralelizaciones, relativa a la
magnitud media del observable.}
\end{table}

\textbf{La energ\'ia concuerda al 0.07\,\% en los cuatro estados.} $C_v$ y $\chi$
dispersan m\'as porque son varianzas estimadas con 100 r\'eplicas: eso es ruido
Monte Carlo, no desacuerdo entre m\'etodos.

\begin{figure}[h]\centering
\includegraphics[width=\textwidth]{figures/obs_points.pdf}
\caption{Observables f\'isicos del estado \texttt{points} (skyrmiones) a lo largo
de la escalera, para las seis paralelizaciones. Filas: cantidad; columnas:
corrida; eje $y$ compartido por fila.}
\end{figure}

\subsection{Pero las realizaciones son distintas}

Comparando las configuraciones finales a \SI{0.1}{\kelvin} entre corridas:

\begin{table}[h]\centering
\begin{tabular}{lrr}
\toprule
estado & correlaci\'on espacial & $Q$ final \\
\midrule
points & $+0.878$ & $-11.26,\,-11.13,\,-11.22,\,-11.20,\,-11.23,\,-11.27$ \\
lines  & $-0.038$ & $-0.33,\,-0.30,\,-0.39,\,-0.25,\,-0.37,\,-0.28$ \\
\bottomrule
\end{tabular}
\end{table}

En \texttt{points} los skyrmiones aparecen en \emph{posiciones distintas} pero las
seis corridas encuentran \textbf{el mismo n\'umero} ($Q$ concuerda al 1.24\,\%).
En \texttt{lines} la correlaci\'on espacial es nula: dos laberintos sin ninguna
relaci\'on geom\'etrica.

\begin{figure}[h]\centering
\includegraphics[width=\textwidth]{figures/trans_points.pdf}
\caption{Transici\'on de estados: $s_z$ de la capa central a lo largo del
recocido, de \SI{20}{\kelvin} a \SI{0.1}{\kelvin}, para \texttt{points}. Las tres
corridas llegan al mismo n\'umero de skyrmiones por caminos distintos.}
\end{figure}

\subsection{Interpretaci\'on}

\textbf{La paralelizaci\'on no afecta la f\'isica; afecta la realizaci\'on
concreta.} Cambiar c\'omo se reparten las cadenas cambia los flujos de n\'umeros
aleatorios: son cadenas de Markov distintas que muestrean la \emph{misma}
distribuci\'on de Boltzmann. Lo que es propiedad de la distribuci\'on ---energ\'ia,
magnetizaci\'on, n\'umero de skyrmiones--- coincide; lo que es una muestra
concreta ---d\'onde cae cada skyrmi\'on--- difiere, y \emph{debe} diferir. Si dos
corridas con semillas distintas dieran la misma imagen, eso s\'i ser\'ia un
defecto.

El estado \texttt{lines} ilustra el punto: con $H_{\mathrm{ex}}=0.022$ el sistema
queda casi degenerado, y existen much\'isimas configuraciones laber\'inticas de
energ\'ia pr\'acticamente igual. Por eso su energ\'ia concuerda al 0.01\,\%
mientras $Q$ dispersa 44\,\%: esos laberintos son energ\'eticamente equivalentes
y geom\'etricamente distintos. Si fuese un defecto de paralelizaci\'on, la
energ\'ia tambi\'en se mover\'ia.

\section{Costos}

H100 a \SI{0.001097}{USD} por segundo (tarifa publicada, verificada).

\begin{table}[h]\centering
\begin{tabular}{lrrrr}
\toprule
corrida & GPU & reloj & GPU-segundos & USD \\
\midrule
r\'eplicas $1\times$ & 1 & 51.7 min & \num{3103} & 3.40 \\
r\'eplicas $2\times$ & 2 & 27.6 min & \num{3307} & 3.63 \\
r\'eplicas $4\times$ & 4 & 16.9 min & \num{4050} & 4.44 \\
ambas $1\times$ & 1 & 62.1 min & \num{3727} & 4.09 \\
ambas $2\times$ & 2 & 30.1 min & \num{3612} & 3.96 \\
ambas $4\times$ & 4 & 16.4 min & \num{3926} & 4.31 \\
\midrule
\textbf{total} & & & \num{21726} & \textbf{23.83} \\
\bottomrule
\end{tabular}
\end{table}

\textbf{El reloj baja $3\times$ mientras el costo sube 30\,\%.} De r\'eplicas
$1\times$ a $4\times$: de 51.7 a 16.9 minutos, de \$3.40 a \$4.44. El 24\,\% que
falta para el escalado ideal se paga en dinero. Si el plazo no aprieta, una sola
GPU es la opci\'on m\'as barata.

Las seis corridas produjeron \num{2400} estados recocidos con observables
completos por \$23.83: \textbf{\$0.0099 por estado}.

\section{L\'imites declarados}

\begin{itemize}
\item \textbf{El coloreado es exacto solo para las capas 1 y 3.} Bajo
$(x+y+z)\bmod 2$ las capas 2 y 4 conectan sitios de la misma paridad, de modo que
con $J_2$ o $J_4$ no nulos la actualizaci\'on simult\'anea deja de ser un
reordenamiento exacto. Las corridas aqu\'i tienen $J_2=J_3=J_4=0$; el 18.9\,\% del
conjunto de datos liberado no, y ese sesgo no est\'a medido.

\item \textbf{Los ritmos por barrido extrapolan \'ordenes de magnitud, no
diferencias del 10--20\,\%.} Al medir las escaleras completas los errores de
extrapolaci\'on resultaron entre $-22.5$\,\% y $+18.7$\,\%, con dos efectos que
empujan en direcciones opuestas: los observables faltantes bajan la estimaci\'on
y el costo fijo por lanzamiento la sube. No existe factor de correcci\'on.

\item \textbf{Extrapolar a lo largo del r\'egimen saturado es exacto; interpolar
hacia dentro del no saturado no.} Una estimaci\'on a 4 cadenas por lanzamiento,
obtenida interpolando entre puntos medidos de 1 y 25, result\'o baja en 48\,\%.

\item \textbf{Las barras de error ajustadas pueden medir la cantidad
equivocada.} Un hallazgo previo de este estudio ---una inversi\'on entre dos
agrupaciones, reportada con 154 y 301 sigma--- result\'o ser un artefacto de
cronometrar el lazo de actualizaci\'on pelado en lugar de una escalera de
producci\'on. Las repeticiones concordaban a \SI{0.0001}{\milli\second} por
barrido, lo que fij\'o el ritmo con precisi\'on pero sobre la cantidad
equivocada.

\item \textbf{Los brazos escalares son CPython.} Las ganancias contra ellos son
frente a una implementaci\'on directa en Python, no frente a c\'odigo secuencial
compilado.
\end{itemize}

\begin{figure}[h]\centering
\includegraphics[width=\textwidth]{figures/extrap_vs_measured.pdf}
\caption{Extrapolaci\'on contra medici\'on para las seis escaleras completas. El
error medio absoluto es 15.2\,\% y el rango va de $-22.5$\,\% a $+18.7$\,\%.}
\end{figure}

\end{document}
